\documentclass[10pt,a4paper]{amsart}
\usepackage{amsmath,amssymb,amsthm,graphicx}
% Retained minimal preamble for the Stacks Project chapter more-algebra.tex (More on Algebra).
% The chapter file pinned-original-source.tex carries only \input{preamble} on its first line, so no
% preamble text lives inside the chapter itself.  The Stacks Project uses one single shared preamble for
% every chapter; the snapshot stages it identically as sources/stacks-chow/preamble.tex and
% sources/stacks-etale-cohomology/preamble.tex.  The block below is copied mechanically, line for line,
% from sources/stacks-chow/preamble.tex lines 14-18 (the Xy-pic setup) and 154-217 (the native theorem
% declarations together with the mathematical macros).  Nothing was added, removed or redefined; the
% wrapper class and the amsmath/amssymb/amsthm/graphicx base packages are supplied by the assembler.
% Chapter: more-algebra.tex -- More on Algebra.
% For commutative diagrams we use Xy-pic
\usepackage[all]{xy}

% We use 2cell for 2-commutative diagrams.
\xyoption{2cell}

% Theorem environments.
%
\theoremstyle{plain}
\newtheorem{theorem}[subsection]{Theorem}
\newtheorem{proposition}[subsection]{Proposition}
\newtheorem{lemma}[subsection]{Lemma}

\theoremstyle{definition}
\newtheorem{definition}[subsection]{Definition}
\newtheorem{example}[subsection]{Example}
\newtheorem{exercise}[subsection]{Exercise}
\newtheorem{situation}[subsection]{Situation}

\theoremstyle{remark}
\newtheorem{remark}[subsection]{Remark}
\newtheorem{remarks}[subsection]{Remarks}

\numberwithin{equation}{subsection}

% Macros
%
\def\lim{\mathop{\mathrm{lim}}\nolimits}
\def\colim{\mathop{\mathrm{colim}}\nolimits}
\def\Spec{\mathop{\mathrm{Spec}}}
\def\Hom{\mathop{\mathrm{Hom}}\nolimits}
\def\Ext{\mathop{\mathrm{Ext}}\nolimits}
\def\SheafHom{\mathop{\mathcal{H}\!\mathit{om}}\nolimits}
\def\SheafExt{\mathop{\mathcal{E}\!\mathit{xt}}\nolimits}
\def\Sch{\mathit{Sch}}
\def\Mor{\mathop{\mathrm{Mor}}\nolimits}
\def\Ob{\mathop{\mathrm{Ob}}\nolimits}
\def\Sh{\mathop{\mathit{Sh}}\nolimits}
\def\NL{\mathop{N\!L}\nolimits}
\def\CH{\mathop{\mathrm{CH}}\nolimits}
\def\proetale{{pro\text{-}\acute{e}tale}}
\def\etale{{\acute{e}tale}}
\def\QCoh{\mathit{QCoh}}
\def\Ker{\mathop{\mathrm{Ker}}}
\def\Im{\mathop{\mathrm{Im}}}
\def\Coker{\mathop{\mathrm{Coker}}}
\def\Coim{\mathop{\mathrm{Coim}}}

% Boxtimes
%
\DeclareMathSymbol{\boxtimes}{\mathbin}{AMSa}{"02}

%
% Macros for moduli stacks/spaces
%
\def\QCohstack{\mathcal{QC}\!\mathit{oh}}
\def\Cohstack{\mathcal{C}\!\mathit{oh}}
\def\Spacesstack{\mathcal{S}\!\mathit{paces}}
\def\Quotfunctor{\mathrm{Quot}}
\def\Hilbfunctor{\mathrm{Hilb}}
\def\Curvesstack{\mathcal{C}\!\mathit{urves}}
\def\Polarizedstack{\mathcal{P}\!\mathit{olarized}}
\def\Complexesstack{\mathcal{C}\!\mathit{omplexes}}
% \Pic is the operator that assigns to X its picard group, usage \Pic(X)
% \Picardstack_{X/B} denotes the Picard stack of X over B
% \Picardfunctor_{X/B} denotes the Picard functor of X over B
\def\Pic{\mathop{\mathrm{Pic}}\nolimits}
\def\Picardstack{\mathcal{P}\!\mathit{ic}}
\def\Picardfunctor{\mathrm{Pic}}
\def\Deformationcategory{\mathcal{D}\!\mathit{ef}}

% --- appended 2026-10-12 (Chief): the project own definitions of the reference/slogan/history
% display environments, copied mechanically from the shared Stacks preamble lines 8-12
% (\usepackage{verbatim} plus the three \newenvironment lines, which the chapter preamble keeps
% verbatim from the project preamble). Nothing added, removed or redefined: these environments
% carry no mathematical content and the project itself typesets them as comments. Without these
% three lines a counted statement whose own text uses them fails with Environment slogan undefined.
\usepackage{verbatim}
\newenvironment{reference}{\comment}{\endcomment}
\newenvironment{slogan}{\comment}{\endcomment}
\newenvironment{history}{\comment}{\endcomment}
\usepackage{cleveref}
\crefname{originalsection}{original section}{original sections}
\Crefname{originalsection}{Original section}{Original sections}
\makeatletter
\newlabel{algebra-section-netherlander}{{00S0}{source}{}{sourcepointer-1}{}}
\newlabel{algebra-section-netherlander@cref}{{[originalsection][0][]00S0}{[1][0][]source}{}{}{}}
\makeatother
\crefname{originallemma}{original lemma}{original lemmas}
\Crefname{originallemma}{Original lemma}{Original lemmas}
\makeatletter
\newlabel{lemma-last-one-flat}{{0653}{source}{}{sourcepointer-2}{}}
\newlabel{lemma-last-one-flat@cref}{{[originallemma][0][]0653}{[1][0][]source}{}{}{}}
\makeatother
\title{Base change of the naive cotangent complex is a quasi-isomorphism over a flat algebra, and derived base change works under a tor amplitude bound}
\author{The Stacks Project Authors}
\date{stacks-project:more-algebra.tex:local-snapshot-20261012; GNU FDL 1.2 or later}
\begin{document}
\maketitle
\setcounter{subsection}{1}
\noindent\textit{Editorial note.} All delivered mathematical text is the author's own, copied from the pinned Stacks Project chapter file more-algebra.tex; no statement or proof step was rewritten, reordered or removed. The counted statement is the labeled lemma environment at source lines 22809-22823 (canonical tag 0FJU) and the counted proof is the proof environment at source lines 22825-22853, which immediately follows it and is its own proof.\par\medskip
\noindent\textit{Editorial note.} The statement displays the cocartesian square of rings in the author's own Xy-pic markup, which is compiled from that markup with the shared preamble, and the naive cotangent complex macro is supplied by the retained preamble.\par\medskip
\noindent\textit{Editorial note.} The preamble span runs to line 90, the current end of the retained preamble, which includes the appended environment definitions; the counted text does not use them, and the span is taken whole so that the file version is unambiguous.\par\medskip
\noindent\textit{Editorial note.} The source URL is the source-level Stacks Project URL, as required when more than one card is produced from one source; the per-claim tag is carried by claim\_key and by this note. The counted statement and proof cite no bibliography entry, so this card needs no bibliography.\par\medskip
\noindent\textit{Allowed standard prerequisites (editorial).}\begin{itemize}
\item A presentation of the algebra by a polynomial algebra over the base gives a two-term complex computing the naive cotangent complex, with the conormal module in degree -1. used at proof lines 22825-22827 and 22845-22851. Fixes the explicit complex whose base change is compared, and expresses the tor amplitude condition through the flatness of the conormal module.
\item For a flat algebra the kernel of the base changed presentation is the tensor product of the original kernel with the new base ring. used at proof lines 22828-22833. Identifies the conormal module of the base changed presentation with the base change of the original one, which gives the isomorphism of complexes.
\item A two-term complex whose degree zero term is free has tor amplitude in the interval from minus one to zero exactly when its degree minus one term is flat, so the derived tensor product coincides with the ordinary one. invoked at proof lines 22849-22852. Turns the ordinary base change isomorphism into a derived one under the tor amplitude hypothesis.
\end{itemize}\medskip
\ifdefined\hypertarget\hypertarget{sourcepointer-1}{}\fi
\noindent\textit{Source reference.} Original section 00S0: sources/stacks-chow/tags.txt, source lines 1022--1022. The reference points to the place where presentations exhibiting the naive cotangent complex are set up, cited at proof lines 22825-22827; it is a pointer to the framework and the computation is inside the counted span.\par\medskip
\ifdefined\hypertarget\hypertarget{sourcepointer-2}{}\fi
\noindent\textit{Source reference.} Original lemma 0653: sources/stacks-more-algebra/pinned-original-source.tex, source lines 16883--16883. The reference converts the tor amplitude hypothesis into flatness of the conormal module, used at proof lines 22849-22850; the conclusion drawn from that flatness is inside the counted span.\par\medskip
\begin{lemma}
\label{lemma-base-change-NL-flat}
Consider a cocartesian diagram of rings
$$
\xymatrix{
B \ar[r] & B' \\
A \ar[r] \ar[u] & A' \ar[u]
}
$$
If $B$ is flat over $A$, then the canonical map
$\NL_{B/A} \otimes_B B' \to \NL_{B'/A'}$ is a quasi-isomorphism.
If in addition $\NL_{B/A}$ has tor-amplitude in $[-1, 0]$
then $\NL_{B/A} \otimes_B^\mathbf{L} B' \to \NL_{B'/A'}$
is a quasi-isomorphism too.
\end{lemma}

\begin{proof}
Choose a presentation $\alpha : P \to B$ as in
Algebra, Section \ref{algebra-section-netherlander}.
Let $I = \Ker(\alpha)$. Set $P' = P \otimes_A A'$ and denote
$\alpha' : P' \to B'$ the corresponding presentation of $B'$ over $A'$.
As $B$ is flat over $A$ we see that $I' = \Ker(\alpha')$ is equal
to $I \otimes_A A'$. Hence
$$
I'/(I')^2 = \Coker(I^2 \otimes_A A' \to I \otimes_A A') =
I/I^2 \otimes_A A' = I/I^2 \otimes_B B'
$$
We have $\Omega_{P'/A'} = \Omega_{P/A} \otimes_A A'$
because both sides have the same basis. It follows that
$\Omega_{P'/A'} \otimes_{P'} B' = \Omega_{P/A} \otimes_P B \otimes_B B'$.
This proves that $\NL(\alpha) \otimes_B B' \to \NL(\alpha')$
is an isomorphism of complexes and hence the first statement holds.

\medskip\noindent
We have
$$
\NL(\alpha) = I/I^2 \longrightarrow \Omega_{P/A} \otimes_P B
$$
as a complex of $B$-modules with $I/I^2$ placed in degree $-1$.
Since the term in degree $0$ is free, this complex has tor-amplitude
in $[-1, 0]$ if and only if $I/I^2$ is a flat $B$-module, see
Lemma \ref{lemma-last-one-flat}.
If this holds, then $\NL(\alpha) \otimes_B^\mathbf{L} B' =
\NL(\alpha) \otimes_B B'$ and we get the second statement.
\end{proof}

\end{document}
